\begin{table}[H]
\centering
\renewcommand{\arraystretch}{1.2}
\begin{tabular}{@{}p{4cm} p{11cm}@{}}
\toprule
\textbf{Use case ID}       & UC-M9-01 \\ \midrule
\textbf{Tên Use case}      & Nhận thông báo về các địa điểm/ người dùng đã theo dõi \\ \midrule
\textbf{Các tác nhân}      & Traveler \\ \midrule
\textbf{Mô tả}             & Traveler nhận thông báo khi có cập nhật mới từ các mục mà người dùng đã theo dõi.\\ \midrule
\textbf{Trigger}    & Người dùng chọn biểu tượng “Thông báo” trên thanh điều hướng.\\ \midrule
\textbf{Tiền điều kiện}    & \begin{tabular}[t]{@{}p{10.5cm}@{}}
1. Người dùng đã đăng nhập vào hệ thống \\
2. Người dùng đã theo dõi ít nhất một đối tượng
\end{tabular} \\ \midrule
\textbf{Hậu điều kiện}     & Hệ thống hiển thị danh sách thông báo về chủ đề người dùng quan tâm \\ \midrule
\textbf{Luồng sự kiện chính} & 
\begin{tabular}[t]{@{}p{10.5cm}@{}}
1. Hệ thống kiểm tra định kỳ hoặc nhận sự kiện mới (bài viết mới, địa điểm cập nhật, người dùng đăng bài).\\
2. Hệ thống tạo nội dung thông báo tương ứng.\\
3. Hệ thống gửi thông báo đến người dùng qua push notification.\\
4. Người dùng xem thông báo trong trung tâm thông báo.
\end{tabular} \\ \midrule
\textbf{Luồng ngoại lệ} &
\begin{tabular}[t]{@{}p{10.5cm}@{}}
4a. Nếu có lỗi kết nối, hệ thống hiển thị thông báo “Lỗi kết nối”\\
\end{tabular} \\
\bottomrule
\end{tabular}
\caption{Mô tả Use Case: Nhận thông báo từ các nội dung theo dõi}
\label{usecase-7-02}
\end{table}

\begin{table}[H]
\centering
\renewcommand{\arraystretch}{1.2}
\begin{tabular}{@{}p{4cm} p{11cm}@{}}
\toprule
\textbf{Use case ID}       & UC-M9-02 \\ \midrule
\textbf{Tên Use case}      & Xem nhật ký hoạt động của cá nhân \\ \midrule
\textbf{Các tác nhân}      & Traveler, Provider \\ \midrule
\textbf{Mô tả}             & Người dùng xem lại lịch sử hoạt động của cá nhân.\\ \midrule
\textbf{Trigger}    & Người dùng chọn biểu tượng “Nhật ký hoạt động”.\\ \midrule
\textbf{Tiền điều kiện}    & Người dùng đã đăng nhập vào hệ thống \\ \midrule
\textbf{Hậu điều kiện}     & Hệ thống hiển thị danh sách hoạt động của người dùng \\ \midrule
\textbf{Luồng sự kiện chính} & 
\begin{tabular}[t]{@{}p{10.5cm}@{}}
1. Hệ thống hiển thị danh sách hoạt động gần nhất theo ngày.\\
2. Người dùng có thể lọc theo loại khoảng thời gian (1 ngày/tuần/tháng trước).
\end{tabular} \\ \midrule
\textbf{Luồng ngoại lệ} &
\begin{tabular}[t]{@{}p{10.5cm}@{}}
1a. Nếu có lỗi kết nối, hệ thống hiển thị thông báo “Lỗi kết nối”\\
\end{tabular} \\
\bottomrule
\end{tabular}
\caption{Mô tả Use Case: Xem nhật ký hoạt động của cá nhân}
\label{usecase-9-02}
\end{table}

\begin{table}[H]
\centering
\renewcommand{\arraystretch}{1.2}
\begin{tabular}{@{}p{4cm} p{11cm}@{}}
\toprule
\textbf{Use case ID}       & UC-M9-03 \\ \midrule
\textbf{Tên Use case}      & Xem nhật ký hoạt động của người dùng \\ \midrule
\textbf{Các tác nhân}      & System admin \\ \midrule
\textbf{Mô tả}             & Quản trị viên xem lịch sử hoạt động của người dùng trong hệ thống\\ \midrule
\textbf{Trigger}    & Quản trị viên chọn trang “Lịch sử hoạt động”.\\ \midrule
\textbf{Tiền điều kiện}    & Người dùng đã đăng nhập vào hệ thống với vai trò là quản trị viên \\ \midrule
\textbf{Hậu điều kiện}     & Hệ thống hiển thị danh sách hoạt động của người dùng \\ \midrule
\textbf{Luồng sự kiện chính} & 
\begin{tabular}[t]{@{}p{10.5cm}@{}}
1. Hệ thống hiển thị danh sách hoạt động gần nhất theo ngày.\\
2. Quản trị viên tiêu chí lọc (theo người dùng, loại hoạt động, ngày, trạng thái…).\\
3. Hệ thống truy xuất và hiển thị kết quả..
\end{tabular} \\ \midrule
\textbf{Luồng ngoại lệ} &
\begin{tabular}[t]{@{}p{10.5cm}@{}}
1a. Nếu có lỗi kết nối, hệ thống hiển thị thông báo “Lỗi kết nối”\\
\end{tabular} \\
\bottomrule
\end{tabular}
\caption{Mô tả Use Case: Xem nhật ký hoạt động của người dùng}
\label{usecase-9-03}
\end{table}

\begin{table}[H]
\centering
\renewcommand{\arraystretch}{1.2}
\begin{tabular}{@{}p{4cm} p{11cm}@{}}
\toprule
\textbf{Use case ID}       & UC-M9-04 \\ \midrule
\textbf{Tên Use case}      & Xem danh sách danh mục của hệ thống \\ \midrule
\textbf{Các tác nhân}      & System admin \\ \midrule
\textbf{Mô tả}             & Quản trị viên xem danh sách các danh mục (địa điểm, chủ đề, loại nội dung) hiện có trong hệ thống.\\ \midrule
\textbf{Trigger}    & Quản trị viên chọn trang “Danh mục hệ thống”.\\ \midrule
\textbf{Tiền điều kiện}    & Người dùng đã đăng nhập vào hệ thống với vai trò là quản trị viên \\ \midrule
\textbf{Hậu điều kiện}     & Hệ thống hiển thị danh sách các nội dung \\ \midrule
\textbf{Luồng sự kiện chính} & 
\begin{tabular}[t]{@{}p{10.5cm}@{}}
1. Hệ thống hiển thị danh sách các nội dung theo từng danh mục (tên, mô tả, trạng thái, ngày cập nhật).\\
2. Quản trị viên có thể lọc hoặc sắp xếp danh sách (theo loại, trạng thái, thời gian).\\
3. Hệ thống truy xuất và hiển thị kết quả.
\end{tabular} \\ \midrule
\textbf{Luồng ngoại lệ} &
\begin{tabular}[t]{@{}p{10.5cm}@{}}
1a. Nếu có lỗi kết nối, hệ thống hiển thị thông báo “Lỗi kết nối”\\
\end{tabular} \\
\bottomrule
\end{tabular}
\caption{Mô tả Use Case: Xem danh sách danh mục của hệ thống}
\label{usecase-9-04}
\end{table}

\begin{table}[H]
\centering
\renewcommand{\arraystretch}{1.2}
\begin{tabular}{@{}p{4cm} p{11cm}@{}}
\toprule
\textbf{Use case ID}       & UC-M9-05 \\ \midrule
\textbf{Tên Use case}      & Thêm, cập nhật, xóa, sửa nội dung của danh mục \\ \midrule
\textbf{Các tác nhân}      & System admin \\ \midrule
\textbf{Mô tả}             & Quản trị viên quản lý nội dung trong các danh mục của hệ thống.\\ \midrule
\textbf{Trigger}    & Quản trị viên chọn biểu tượng ứng với thao tác: thêm, sửa, hoặc xóa nội dung theo từng danh mục.
\\ \midrule
\textbf{Tiền điều kiện}    & Người dùng đã đăng nhập vào hệ thống với vai trò là quản trị viên \\ \midrule
\textbf{Hậu điều kiện}     & Hệ thống hiển thị danh sách các nội dung đã được cập nhật. \\ \midrule
\textbf{Luồng sự kiện chính} & 
\begin{tabular}[t]{@{}p{10.5cm}@{}}
1. Hệ thống hiển thị form nhập thông tin hoặc xác nhận thao tác.
2. Quản trị viên nhập dữ liệu hợp lệ.
3. Hệ thống kiểm tra dữ liệu và lưu thay đổi.
4. Thông báo kết quả thao tác.
\end{tabular} \\ \midrule
\textbf{Luồng ngoại lệ} &
\begin{tabular}[t]{@{}p{10.5cm}@{}}
1a. Nếu có lỗi kết nối, hệ thống hiển thị thông báo “Lỗi kết nối”\\
\end{tabular} \\
\bottomrule
\end{tabular}
\caption{Mô tả Use Case: Quản lý nội dung của danh mục}
\label{usecase-9-04}
\end{table}
